%% kapa: exam=1 round=33 year=2022 period=2 form=A subject=appraisal_law subject_ko=감정평가관계법규 seq=18 number=18
%% source: Q-Net 감정평가사 기출문제 (cst00309) artlSeq=5209997 file=raw/1008/5209997/2022년도 제33회 감정평가사 1차시험 2교시(감정평가 관계법규, 회계학).pdf
%% conversion: kordoc 4.24.51 -> tools/gen1.py (auto); review=proofread-sonnet; flags=none
\begin{question}{18}
부동산 가격공시에 관한 법령상 개별공시지가에 관한 설명으로 옳지 않은 것은?
\begin{choices}
  \item 표준지로 선정된 토지에 대하여는 개별공시지가를 결정·공시하지 아니할 수 있다.
  \item 개별토지 가격 산정의 타당성에 대한 감정평가법인등의 검증을 생략하려는 경우 개발사업이 시행되는 토지는 검증 생략 대상 토지로 선정해서는 안 된다.
  \item 개별토지 가격 산정의 타당성 검증을 의뢰할 감정평가법인등을 선정할 때 선정기준일부터 직전 1년간 과태료처분을 2회 받은 감정평가법인등은 선정에서 배제된다.
  \item 개별공시지가 조사·산정의 기준에는 토지가격비준표의 사용에 관한 사항이 포함되어야 한다.
  \item 개별공시지가에 이의가 있는 자는 그 결정·공시일부터 30일 이내에 서면으로 시장·군수 또는 구청장에게 이의를 신청할 수 있다.
\end{choices}
\end{question}
\answer{3}
